\documentclass{ieeeaccess}
\usepackage{cite}
\usepackage{amsmath,amssymb,amsfonts}
\usepackage{algorithmic}
\usepackage{algorithm}
\usepackage{graphicx}
\usepackage{textcomp}
\usepackage{amsthm}
\usepackage{pifont}
\usepackage{hyperref}
\hypersetup{
    colorlinks=true,
    linkcolor=cyan,
    filecolor=magenta,      
    urlcolor=cyan,
    citecolor=cyan,
    bookmarks=true,
    bookmarksopen=true,
    bookmarksdepth=3,
}
\usepackage{booktabs}
\usepackage{multirow}

\theoremstyle{definition}
\newtheorem{definition}{Definition}

\usepackage{bm}
\makeatletter
\AtBeginDocument{\DeclareMathVersion{bold}
\SetSymbolFont{operators}{bold}{T1}{times}{b}{n}
\SetSymbolFont{NewLetters}{bold}{T1}{times}{b}{it}
\SetMathAlphabet{\mathrm}{bold}{T1}{times}{b}{n}
\SetMathAlphabet{\mathit}{bold}{T1}{times}{b}{it}
\SetMathAlphabet{\mathbf}{bold}{T1}{times}{b}{n}
\SetMathAlphabet{\mathtt}{bold}{OT1}{pcr}{b}{n}
\SetSymbolFont{symbols}{bold}{OMS}{cmsy}{b}{n}
\renewcommand\boldmath{\@nomath\boldmath\mathversion{bold}}}
\makeatother

\def\BibTeX{{\rm B\kern-.05em{\sc i\kern-.025em b}\kern-.08em
    T\kern-.1667em\lower.7ex\hbox{E}\kern-.125emX}}

%Your document starts from here ___________________________________________________
\begin{document}
\history{Date of publication xxxx 00, 0000, date of current version xxxx 00, 0000.}
\doi{10.1109/ACCESS.2024.0429000}

\title{DPFed-GridGuard: An Enhanced Federated Learning Framework with Multi-Dimensional Differential Privacy for Intrusion Detection in Decentralized Smart Grids}
\author{\uppercase{Anass BETOUIL}\authorrefmark{1},
\uppercase{Samia EL HADDOUTI\authorrefmark{1,2,3}, and Habiba CHAOUI}.\authorrefmark{1}}
\address[1]{Laboratory of Advanced Systems Engineering, National School of Applied Sciences, Ibn Tofail University, Kenitra, Morocco}
\address[2]{National Higher School For Computer Science and Systems Analysis, Mohamed V University, Rabat, Morocco}
\address[3]{National Center for Scientific and Technical Research, Rabat, Morocco}

\markboth
{BETOUIL \headeretal: DPFed-GridGuard}
{BETOUIL \headeretal: DPFed-GridGuard}

\corresp{Corresponding author: Anass BETOUIL (e-mail: anass.betouil@uit.ac.ma).}


\begin{abstract}
The widespread deployment of smart grid systems, generating vast amounts of sensitive data, creates numerous vulnerabilities and makes these systems prime targets for sophisticated cyber attacks. This environment requires advanced anomaly detection methods that can ensure high detection accuracy while rigorously preserving data privacy. Federated Learning (FL) has emerged a promising approach for training models collaboratively across distributed systems without revealing raw data. Yet, the deployment of FL in decentralized smart grid systems presents several challenges, particularly due to heterogeneous data distributions, and the risk that model updates may leak sensitive information through inference or poisoning attacks. To address these issues, we present DPFed-GridGuard, a novel FL framework enhanced with multi-dimensional differential privacy techniques to secure sensitive consumption data while maintaining high detection performance. Specifically, DPFed-GridGuard integrates four complementary mechanisms: (1) Context-Aware Local Differential Privacy (CA-LDP) that adapts noise to feature sensitivity, (2) Cluster-Adaptive Differential Privacy (CADP) that customizes privacy levels for heterogeneous clusters, (3) Selective Homomorphic Encryption (S-HE) which encrypts critical model components to reduce overhead, and (4) a Utility-Aware Noise Scheduler (UANS) that optimizes privacy budgets dynamically across training rounds. Experiment evaluation on three diverse real-world datasets demonstrate that DPFed-GridGuard achieves superior anomaly detection accuracy up to y\% with utility retention exceeding X\%, while robustly protecting privacy and minimizing communication and computational costs. Accordingly, the proposed framework offers a practical and scalable solution for secure and privacy-preserving anomaly detection in decentralized smart grid environments.
\end{abstract}

\begin{keywords}
Federated learning, smart grid security, differential privacy, anomaly detection, privacy-preserving
\end{keywords}

\titlepgskip=-21pt

\maketitle

\section{Introduction}
\label{sec:introduction}
\PARstart{T}{his} transformation of traditional power grids into smart grids represents one of the most significant infrastructure upgrades of the 21st century. By integrating advanced sensing technologies, bidirectional communication systems, and intelligent control mechanisms, smart grids enable unprecedented optimization and control capabilities \cite{farhangi2010path, gungor2011smart}. However, the granular data that enables efficient grid operation also exposes sensitive information about consumers and critical infrastructure. Smart meter readings can reveal detailed household activities and industrial operations \cite{molina2010private, mckenna2012smart}, while infrastructure measurements could be exploited to identify system vulnerabilities.

The evolving cybersecurity landscape demands robust anomaly detection systems, yet traditional centralized approaches require data aggregation that creates privacy risks and faces regulatory scrutiny under frameworks like GDPR and CCPA \cite{voigt2017gdpr, pardau2018california}. While federated learning offers a promising alternative by keeping data local \cite{mcmahan2017communication, kairouz2021advances}, recent research shows that standard federated learning implementations remain vulnerable to privacy attacks, as gradient updates can still leak sensitive information about training data \cite{jithish2023comprehensive, shrestha2024anomaly, mothukuri2021survey}. Smart grid environments present unique challenges with their heterogeneous data distributions across regions, customer types, and time periods, requiring specialized approaches for effective anomaly detection \cite{wen2021feddetect, perez2024review}.

Existing solutions fail to address the critical requirements identified through industry consultation: maintaining operational effectiveness ($\geq90$\% accuracy), providing formal privacy guarantees as required by regulations, ensuring computational efficiency on resource-constrained devices, and handling inherent data heterogeneity. To address these challenges, we present DPFed-GridGuard, a comprehensive privacy-enhanced federated learning framework specifically designed for smart grid anomaly detection.

Our framework makes six key contributions: (1) An integrated privacy framework combining four complementary techniques designed for practical deployment; (2) Context-Aware Local Differential Privacy (CA-LDP) that applies targeted noise based on feature sensitivity, improving utility over uniform approaches; (3) Cluster-Adaptive Differential Privacy (CADP) that customizes privacy for heterogeneous deployments, maintaining high performance under non-IID conditions; (4) Selective Homomorphic Encryption (S-HE) that encrypts only critical model dimensions, reducing computational overhead; (5) Utility-Aware Noise Scheduler (UANS) that dynamically allocates privacy budgets, accelerating convergence; and (6) Comprehensive evaluation on three real-world datasets demonstrating exceptional performance retention while providing formal privacy guarantees.

The remainder of this paper is organized as follows: Section~\ref{sec:related_work} reviews related work, Section~\ref{sec:system_model} presents the system model, Section~\ref{sec:framework} details the DPFed-GridGuard framework, Section~\ref{sec:experiments} describes the experimental setup, Section~\ref{sec:results} presents results, Section~\ref{sec:discussion} discusses implications and limitations, and Section~\ref{sec:conclusion} concludes the paper.

\section{Related Work}
\label{sec:related_work}

The intersection of federated learning, privacy preservation, and smart grid security has attracted significant research attention in recent years. This section reviews the most relevant work in these areas, highlighting both the progress made and the gaps that our work addresses.

\subsection{Federated Learning for Critical Infrastructure}

The application of federated learning to critical infrastructure systems has gained momentum as organizations seek to leverage collective intelligence while maintaining data sovereignty. Recent comprehensive surveys \cite{zhang2024federated, cheng2022review} have highlighted the growing interest in applying federated learning to smart grid applications, while identifying key challenges in privacy preservation and heterogeneous data handling.

One of the recent efforts in this space was the comprehensive survey by Zhang et al. \cite{zhang2024federated}, which applied federated learning to anomaly detection in smart grids. They demonstrated that FL models could achieve comparable anomaly detection performance to centralized approaches while keeping data distributed. However, their work focused primarily on binary classification (normal vs. anomaly) and did not incorporate differential privacy mechanisms, leaving the system vulnerable to gradient-based privacy attacks. Similar limitations were identified in a federated learning framework proposed by Liu et al. \cite{liu2021federated}, which focused on securing power traces but lacked formal privacy guarantees.

Building on the privacy concerns, Shrestha et al. \cite{shrestha2024anomaly} introduced an LSTM-autoencoder based approach for anomaly detection using federated learning in smart electric grids. While their work showed promise in detecting anomalies using sequential patterns, they acknowledged that standard federated learning still reveals information through model updates and suggested the need for additional privacy mechanisms. This need for enhanced privacy has been further emphasized in recent surveys on federated learning security and privacy \cite{cao2024survey}.

Wen et al. \cite{wen2021feddetect} proposed FedDetect, a privacy-preserving federated learning framework specifically for energy theft detection. Their approach incorporated basic differential privacy but focused solely on binary classification (theft vs. normal) and reported accuracy drops of 12-15\% when privacy mechanisms were enabled. This significant utility loss highlights the challenge of balancing privacy and performance in real-world deployments.

A more recent advancement came with the work of Alshehri et al. \cite{alshehri2024deep}, which explored deep anomaly detection for zero-day cyberattacks in electricity theft scenarios. While innovative in targeting unknown attacks, their federated approach still lacked formal privacy guarantees and was limited to theft detection rather than comprehensive anomaly detection across multiple attack types. Recent advances in decentralized approaches \cite{husnoo2024decentralized} have shown promise in addressing some of these limitations through P2P gossip protocols.

\subsection{Privacy-Preserving Federated Learning}

The privacy vulnerabilities of standard federated learning have been well-documented in recent literature. While the foundational work by McMahan et al. \cite{mcmahan2017communication} introduced federated averaging as a way to train models without centralizing data, subsequent research has shown that model updates can still leak significant information \cite{li2020federated}.

Differential privacy has emerged as the most widely adopted approach for providing formal privacy guarantees in federated learning \cite{dwork2014algorithmic}. Ren et al. \cite{ren2023efeddsa} introduced EFedDSA, an efficient differential privacy-based horizontal federated learning approach for smart grid dynamic security assessment. They achieved reasonable privacy-utility trade-offs but focused on security assessment rather than anomaly detection and were limited to specific types of grid measurements.

The challenge of applying differential privacy effectively in federated settings has been explored by several researchers. Zhao et al. \cite{zhao2021differential} proposed a differential privacy-enhanced federated learning method for short-term household load forecasting, demonstrating that careful noise calibration could maintain forecast accuracy. However, their approach used uniform noise addition across all features, missing opportunities for optimization based on feature sensitivity. Recent work by Nguyen et al. \cite{nguyen2024enhancing} has shown improved results using LSTM and GRU architectures with differential privacy for PV feed-in power forecasting.

Recent work has begun exploring adaptive privacy mechanisms. Badar et al. \cite{badar2023privacy} proposed a privacy-preserving energy prediction scheme that adjusted privacy levels based on data characteristics, but their work focused on prediction rather than anomaly detection and did not address the multi-class classification challenges present in comprehensive security monitoring. The integration of differential privacy with clustering approaches has shown promise, as demonstrated by Zhang et al. \cite{zhang2025dpck} in their DPCK framework for smart meter data analysis.

Beyond differential privacy, researchers have explored other privacy-preserving techniques. Ashraf et al. \cite{ashraf2022feddp} proposed FedDP, combining federated learning with additional privacy protections for theft detection. While they achieved good detection rates, their system was designed for binary classification and did not scale to multi-class scenarios. Recent advances in homomorphic encryption \cite{choi2024fednic, islam2025smart} and secure multi-party computation \cite{liu2024privacy} offer complementary approaches to enhance privacy in federated learning settings.

\subsection{State of the Art}

Table~\ref{tab:sota_comparison} provides a detailed comparison of DPFed-GridGuard with existing approaches, highlighting the unique contributions of our work. Recent developments include risk assessment frameworks \cite{song2024data}, heterogeneous federated learning approaches \cite{zhao2024heterofl}, and edge intelligence integration \cite{chen2024introducing}.

\begin{table*}[t]
\caption{Comprehensive Comparison with State-of-the-Art Federated Learning Methods for Smart Grids}
\label{tab:sota_comparison}
\centering
\begin{tabular}{lccccccc}
\toprule
Framework & Domain & Max Classes & Privacy Mechanism & Avg. Retention & Non-IID Support & Convergence & Key Innovation \\
\midrule
FedDetect \cite{wen2021feddetect} & Power & Binary & Basic DP & 85-88\% & No & 20 rounds & Energy theft focus \\
EFedDSA \cite{ren2023efeddsa} & Power & 5 & Uniform DP & 87\% & Yes & 25 rounds & Security assessment \\
Li et al. \cite{li2025clustered} & Power & Multi-class & Clustered FL & 86\% & Yes & 18 rounds & FDIA detection \\
Zhang et al. \cite{zhang2025dpck} & Power & Binary & Adaptive DP & 88.7\% & Limited & 22 rounds & Smart meter privacy \\
\textbf{DPFed-GridGuard} & \textbf{Power} & \textbf{Multi-class} & \textbf{CA-LDP + CADP} & \textbf{98.9\%} & \textbf{Yes} & \textbf{20 rounds} & \textbf{Integrated privacy} \\
\bottomrule
\end{tabular}
\end{table*}

\begin{itemize}
    \item Comprehensive Privacy Protection: While existing works either ignore privacy or apply basic mechanisms, DPFed-GridGuard integrates four complementary privacy techniques designed specifically for smart grid characteristics. This integrated approach enables us to achieve superior privacy-utility trade-offs.
    \item Multi-Class Capability: Most existing works focus on binary classification (normal vs. anomalous), which oversimplifies the reality of smart grid security where different types of attacks require different responses. Our framework handles multiple distinct attack classes, providing fine-grained detection capabilities essential for effective incident response.
    \item Heterogeneity Handling: Real smart grid deployments exhibit significant data heterogeneity due to geographic, demographic, and infrastructure variations. While most existing approaches assume IID data or provide limited non-IID support, our CADP mechanism is specifically designed to handle the heterogeneous nature of smart grid data.
    \item Practical Efficiency: The significant reduction in communication rounds achieved by DPFed-GridGuard is crucial for real-world deployment where communication bandwidth is often limited and expensive. This efficiency gain, combined with our selective encryption approach, makes the framework practical for deployment on existing infrastructure.
    \item Formal Guarantees: Unlike approaches that provide only empirical privacy protection, DPFed-GridGuard offers formal differential privacy guarantees that can be verified and audited. This is essential for regulatory compliance and building trust among participating organizations.
\end{itemize}

These advantages position DPFed-GridGuard as a comprehensive solution for privacy-preserving federated learning in smart grid security applications, addressing both theoretical and practical challenges that have limited the deployment of previous approaches.

\section{System Model and Problem Formulation}
\label{sec:system_model}

In this section, we present the system model for federated learning in smart grid environments and formally define the privacy-preserving anomaly detection problem. We begin by describing the architectural components and their interactions, then formalize the optimization objectives and constraints.

\subsection{Smart Grid Federated Learning Architecture}

Modern smart grids operate as complex, hierarchical systems with multiple layers of data collection, processing, and control. Our federated learning framework leverages this existing infrastructure to enable collaborative anomaly detection while preserving data privacy. Figure~\ref{fig:federated_architecture} illustrates the three-tier architecture we consider, which aligns with recent architectural proposals for edge-cloud collaborative systems \cite{bekara2022federated, asad2022federated}.

\begin{figure*}[h]
\centering
\includegraphics[width=\textwidth]{figures/fl_architecture.pdf}
\caption{Three-tier federated learning architecture for smart grid anomaly detection. The cloud layer hosts the Utility Control Center for global coordination, the fog layer contains substations for regional aggregation, and the edge layer consists of smart meters and PMUs performing local computation.}
\label{fig:federated_architecture}
\end{figure*}

At the \textit{edge layer}, smart meters and PMUs continuously collect high-frequency measurements including voltage magnitudes, phase angles, current flows, and power quality indicators. These devices possess limited computational capabilities but are sufficient for basic model training operations, as demonstrated in recent edge intelligence research \cite{chen2024introducing, anand2024enhancing}. Each edge device maintains local storage for recent measurements and participates in federated learning by training local models on this data.

The \textit{fog layer} consists of substations and regional control centers that serve as intermediate aggregation points. These facilities have more substantial computational resources and act as regional coordinators, aggregating updates from multiple edge devices before forwarding them to the cloud layer. This hierarchical aggregation reduces communication overhead and provides a natural boundary for applying differential privacy mechanisms.

At the \textit{cloud layer}, the Utility Control Center coordinates the overall federated learning process. It maintains the global model, orchestrates training rounds, and distributes updated models back through the hierarchy. Importantly, the cloud layer never receives raw measurement data, only aggregated model updates that have been protected by our privacy mechanisms, following secure aggregation principles \cite{abdellatif2024sdcl}.

\subsection{Formal Problem Definition}

Let us consider a smart grid system with $N$ participating entities (smart meters, substations, and control centers), denoted as $\mathcal{C} = \{C_1, C_2, ..., C_N\}$. Each entity $C_i$ possesses a local dataset $\mathcal{D}_i = \{(x_j, y_j)\}_{j=1}^{n_i}$, where $x_j \in \mathbb{R}^d$ represents a $d$-dimensional feature vector of grid measurements and $y_j \in \{0, 1, ..., K-1\}$ indicates the operational state (normal or one of $K-1$ types of anomalies).

The objective of federated learning in this context is to collaboratively train a global model $w^*$ that minimizes the overall empirical risk:

\begin{equation}
w^* = \arg\min_{w} F(w) = \arg\min_{w} \sum_{i=1}^{N} \frac{n_i}{n} F_i(w)
\end{equation}

where $n = \sum_{i=1}^{N} n_i$ represents the total number of samples across all participants, and $F_i(w)$ denotes the local objective function for participant $i$:

\begin{equation}
F_i(w) = \frac{1}{n_i} \sum_{j=1}^{n_i} \ell(f(x_j; w), y_j)
\end{equation}

Here, $f(\cdot; w)$ represents the anomaly detection model parameterized by $w$, and $\ell(\cdot, \cdot)$ is the loss function (typically cross-entropy for multi-class classification).

\subsection{Privacy Requirements and Threat Model}

Smart grid data contains highly sensitive information that must be protected throughout the federated learning process. We consider the following privacy requirements:

\textbf{Requirement 1 (Data Confidentiality):} Raw measurement data $\mathcal{D}_i$ must never leave the local entity $C_i$. Only processed model updates should be communicated.

\textbf{Requirement 2 (Differential Privacy):} The federated learning process must satisfy $(\varepsilon, \delta)$-differential privacy, formally defined as:

\begin{definition}[Differential Privacy]
A randomized mechanism $\mathcal{M}$ satisfies $(\varepsilon, \delta)$-differential privacy if for any two adjacent datasets $\mathcal{D}$ and $\mathcal{D}'$ differing in at most one record, and for any subset of outputs $S \subseteq \text{Range}(\mathcal{M})$:
\begin{equation}
\Pr[\mathcal{M}(\mathcal{D}) \in S] \leq e^{\varepsilon} \cdot \Pr[\mathcal{M}(\mathcal{D}') \in S] + \delta
\end{equation}
\end{definition}

\textbf{Requirement 3 (Operational Effectiveness):} Despite privacy protections, the system must maintain a minimum detection accuracy of 90\% to meet operational requirements for critical infrastructure protection.

We adopt an honest-but-curious threat model, which reflects realistic scenarios in smart grid deployments where participants follow the protocol correctly but may attempt to infer sensitive information from observed communications. Specifically, we consider the following potential adversaries:

\begin{itemize}
\item \textbf{Curious Aggregator:} The central server honestly performs aggregation but attempts to infer information about individual participants from model updates.
\item \textbf{Curious Participants:} Other participating entities that honestly contribute to training but try to learn about peers' data from shared model parameters.
\item \textbf{External Observers:} Adversaries who can observe communication channels and attempt to extract information from intercepted model updates.
\end{itemize}

Table~\ref{tab:threats} summarizes the specific privacy threats we address and their potential impact on smart grid operations.

\begin{table}[htbp]
\caption{Privacy Threats in Federated Smart Grid Systems}
\label{tab:threats}
\centering
\begin{tabular}{p{2.5cm}p{5cm}}
\toprule
Attack Type & Description and Impact \\
\midrule
Gradient Inference & Adversary analyzes gradient updates to reconstruct training samples, potentially revealing consumption patterns or operational states \\
\midrule
Membership Inference & Determines whether specific measurements were used in training, potentially exposing which customers or facilities participated \\
\midrule
Property Inference & Extracts statistical properties of local datasets, revealing aggregate consumption patterns or attack frequencies \\
\midrule
Model Inversion & Reconstructs representative training samples from model parameters, potentially exposing typical operational patterns \\
\bottomrule
\end{tabular}
\end{table}

\subsection{Data Heterogeneity in Smart Grids}

One of the fundamental challenges in applying federated learning to smart grids is the inherent heterogeneity of data across different participants \cite{li2025federated, samek2025intelligent}. This heterogeneity manifests in several ways:

\textbf{Spatial Heterogeneity:} Different geographic regions exhibit distinct consumption patterns based on climate, demographics, and industrial composition. For example, residential areas show pronounced morning and evening peaks, while industrial zones maintain relatively constant consumption during working hours.

\textbf{Temporal Heterogeneity:} Consumption patterns vary significantly across different time scales—daily, weekly, and seasonal variations create non-stationary distributions that challenge traditional federated learning assumptions.

\textbf{Attack Pattern Heterogeneity:} Different regions face different types of cyber threats based on their infrastructure importance and vulnerability. Critical substations may experience more sophisticated attacks, while residential areas might see more opportunistic intrusions.

\textbf{Infrastructure Heterogeneity:} The age, type, and configuration of equipment varies across the grid, leading to differences in measurement quality, sampling rates, and available features.

To quantify this heterogeneity, we use the following metrics:

\begin{definition}[Data Distribution Divergence]
For two participants $i$ and $j$ with label distributions $p_i$ and $p_j$, the distribution divergence is measured using KL divergence:
\begin{equation}
D_{KL}(p_i || p_j) = \sum_{k=0}^{K-1} p_i(k) \log \frac{p_i(k)}{p_j(k)}
\end{equation}
\end{definition}

\begin{definition}[System Heterogeneity Level]
The overall system heterogeneity $H$ is computed as the average pairwise divergence:
\begin{equation}
H = \frac{2}{N(N-1)} \sum_{i=1}^{N-1} \sum_{j=i+1}^{N} D_{KL}(p_i || p_j)
\end{equation}
\end{definition}

Table~\ref{tab:heterogeneity} shows typical heterogeneity levels and their impact on federated learning performance.

\begin{table}[htbp]
\caption{Data Heterogeneity Scenarios and Performance Impact}
\label{tab:heterogeneity}
\centering
\begin{tabular}{lccc}
\toprule
Scenario & H Value & Distribution Type & Expected Gap \\
\midrule
IID & < 0.1 & Uniform & 0-5\% \\
Mild Non-IID & 0.1-0.5 & Moderately Skewed & 5-10\% \\
Strong Non-IID & > 0.5 & Highly Skewed & 10-20\% \\
\bottomrule
\end{tabular}
\end{table}

\subsection{Privacy-Constrained Optimization Problem}

Given these requirements and challenges, we formulate the privacy-preserving federated learning problem for smart grids as a constrained optimization problem:

\begin{align}
\min_{w, \mathcal{M}} \quad & F(w) = \sum_{i=1}^{N} \frac{n_i}{n} F_i(w) \\
\text{subject to} \quad & \mathcal{M} \text{ satisfies } (\varepsilon, \delta)\text{-DP} \\
& \text{Accuracy}(w) \geq 0.90 \\
& \text{Rounds} \leq R_{\max} \\
& \text{Compute}_{i} \leq C_{\max} \quad \forall i \in \{1, ..., N\}
\end{align}

where $\mathcal{M}$ represents the privacy mechanism, $R_{\max}$ denotes the maximum number of communication rounds, and $C_{\max}$ represents the computational budget for each participant.

This formulation captures the fundamental trade-offs in privacy-preserving federated learning for smart grids: we seek to minimize the global loss (maximize anomaly detection performance) while satisfying strict privacy constraints, maintaining operational accuracy requirements, and respecting practical limitations on communication and computation.

\section{DPFed-GridGuard Framework}
\label{sec:framework}

This section presents the detailed design of DPFed-GridGuard, our privacy-enhanced federated learning framework for smart grid anomaly detection. We describe each of the four core privacy mechanisms and explain how they work together to achieve superior privacy-utility trade-offs.

\subsection{Framework Overview and Design Philosophy}

DPFed-GridGuard is built on four fundamental design principles that guide all architectural decisions:

\textbf{Principle 1: Context-Aware Privacy.} Rather than applying uniform privacy protection across all data and model components, we recognize that different aspects of smart grid data have varying sensitivity levels. This principle drives our development of mechanisms that adapt privacy protection based on context.

\textbf{Principle 2: Hierarchical Aggregation.} We leverage the natural hierarchy of smart grid infrastructure to implement privacy mechanisms at appropriate levels, reducing overhead while maintaining strong guarantees.

\textbf{Principle 3: Adaptive Resource Allocation.} Both privacy budget and computational resources are precious in federated learning. Our framework dynamically allocates these resources where they provide the most benefit.

\textbf{Principle 4: Practical Deployability.} Every design decision considers real-world constraints, ensuring that the framework can be deployed on existing infrastructure without requiring significant upgrades.

The overall workflow of DPFed-GridGuard follows an enhanced federated learning protocol with integrated privacy mechanisms at each stage. At the beginning of each training round, the central server broadcasts the current global model to participating clients. Each client then applies our privacy mechanisms during local training and before uploading updates. The server aggregates these privacy-protected updates to form the new global model, continuing this process until convergence.

\subsection{Context-Aware Local Differential Privacy (CA-LDP)}

Traditional differential privacy mechanisms add uniform noise across all features, which can unnecessarily degrade model utility. Our Context-Aware Local Differential Privacy mechanism recognizes that not all features in smart grid data are equally sensitive or equally important for anomaly detection.

\subsubsection{Feature Sensitivity Analysis}

Before applying differential privacy, CA-LDP analyzes the sensitivity of each feature in the smart grid data. We define feature sensitivity based on two factors:

\begin{equation}
s_i = \alpha \cdot \text{Privacy}_i + (1-\alpha) \cdot \text{Utility}_i
\end{equation}

where $\text{Privacy}_i$ measures how much feature $i$ reveals about individual consumption patterns, and $\text{Utility}_i$ captures its importance for anomaly detection. The parameter $\alpha \in [0,1]$ allows operators to balance these concerns based on their specific requirements.

To compute privacy sensitivity, we analyze the mutual information between each feature and potentially sensitive attributes:

\begin{equation}
\text{Privacy}_i = \frac{I(X_i; S)}{\max_j I(X_j; S)}
\end{equation}

where $I(X_i; S)$ represents the mutual information between feature $X_i$ and sensitive attributes $S$ (such as individual consumption levels or appliance usage patterns).

For utility importance, we use a combination of feature correlation with labels and contribution to model performance:

\begin{equation}
\text{Utility}_i = \beta \cdot |\text{Corr}(X_i, Y)| + (1-\beta) \cdot \text{SHAP}_i
\end{equation}

where $\text{SHAP}_i$ represents the average SHAP (SHapley Additive exPlanations) value for feature $i$ from a baseline model \cite{lundberg2017shap, lv2022review, ali2024enhancing}.

\subsubsection{Adaptive Noise Calibration}

Once we have computed sensitivity scores, CA-LDP applies differential privacy with adaptive noise calibration:

\begin{algorithm}
\caption{Context-Aware Local Differential Privacy}
\label{alg:ca_ldp}
\begin{algorithmic}[1]
\REQUIRE Local data $\mathcal{D}$, sensitivity scores $S = \{s_1, ..., s_d\}$, privacy budget $\varepsilon$
\ENSURE Privacy-protected gradient $\tilde{\nabla}F$
\STATE Train local model on $\mathcal{D}$ to compute gradient $\nabla F$
\STATE Compute feature-specific noise scales:
\FOR{each feature dimension $i \in \{1, ..., d\}$}
    \STATE $\lambda_i \leftarrow \frac{s_i \cdot \Delta_i}{\varepsilon}$ where $\Delta_i$ is the sensitivity bound
\ENDFOR
\STATE Generate adaptive noise:
\FOR{each gradient component $g_i$}
    \STATE $\eta_i \sim \text{Laplace}(0, \lambda_i)$
    \STATE $\tilde{g}_i \leftarrow g_i + \eta_i$
\ENDFOR
\STATE Apply gradient clipping: $\tilde{\nabla}F \leftarrow \text{clip}(\tilde{g}, C)$
\RETURN $\tilde{\nabla}F$
\end{algorithmic}
\end{algorithm}

The key innovation here is that features identified as highly sensitive (high privacy risk) receive more noise, while features crucial for anomaly detection but low in privacy risk receive less noise. This targeted approach maintains the overall privacy guarantee while improving utility.

\subsection{Cluster-Adaptive Differential Privacy (CADP)}

Smart grid environments exhibit significant heterogeneity where residential areas have different privacy needs than industrial zones, and regions with frequent attacks may need different protection levels than stable areas. CADP addresses this by clustering clients based on their characteristics and assigning appropriate privacy parameters.

\subsubsection{Client Clustering}

We characterize each client using a feature vector that captures both data properties and operational context:

\begin{equation}
\phi_i = [\text{size}_i, \text{entropy}_i, \text{attack\_rate}_i, \text{type}_i, \text{variance}_i]
\end{equation}

where:
- $\text{size}_i$ is the number of local training samples
- $\text{entropy}_i$ measures the class distribution entropy
- $\text{attack\_rate}_i$ is the historical frequency of anomalies
- $\text{type}_i$ encodes the client type (residential, commercial, industrial)
- $\text{variance}_i$ captures the variability in measurements

We then apply k-means clustering to group similar clients:

\begin{algorithm}
\caption{Cluster-Adaptive Privacy Assignment}
\label{alg:cadp}
\begin{algorithmic}[1]
\REQUIRE Client features $\{\phi_1, ..., \phi_N\}$, base privacy budget $\varepsilon_0$, number of clusters $K$
\ENSURE Cluster assignments $\{c_1, ..., c_N\}$ and privacy budgets $\{\varepsilon_1, ..., \varepsilon_K\}$
\STATE Apply k-means clustering: $\mathcal{C}_1, ..., \mathcal{C}_K \leftarrow \text{KMeans}(\{\phi_i\}, K)$
\FOR{each cluster $k \in \{1, ..., K\}$}
    \STATE Compute cluster statistics:
    \STATE $\quad \bar{s}_k \leftarrow$ average data size in cluster $k$
    \STATE $\quad \bar{a}_k \leftarrow$ average attack rate in cluster $k$
    \STATE $\quad \bar{v}_k \leftarrow$ average variance in cluster $k$
    \STATE Determine adaptive privacy budget:
    \STATE $\quad \varepsilon_k \leftarrow \varepsilon_0 \cdot f(\bar{s}_k, \bar{a}_k, \bar{v}_k)$
\ENDFOR
\STATE Assign each client $i$ to cluster $c_i$ with privacy budget $\varepsilon_{c_i}$
\RETURN Cluster assignments and privacy budgets
\end{algorithmic}
\end{algorithm}

\subsubsection{Privacy Budget Adaptation}

The adaptation function $f$ considers multiple factors to determine appropriate privacy levels:

\begin{equation}
f(\bar{s}, \bar{a}, \bar{v}) = \frac{1}{Z} \left( w_s \cdot \frac{\bar{s}}{s_{\max}} + w_a \cdot (1 - \bar{a}) + w_v \cdot \frac{1}{1 + \bar{v}} \right)
\end{equation}

where $Z$ is a normalization constant ensuring the average privacy budget remains $\varepsilon_0$, and $w_s$, $w_a$, $w_v$ are weights that can be tuned based on deployment priorities.

This approach ensures that:
- Clients with more data (higher $\bar{s}$) can afford slightly less aggressive privacy protection
- Clients in high-attack areas (higher $\bar{a}$) receive stronger privacy protection
- Clients with stable patterns (lower $\bar{v}$) can use less noise without compromising privacy

\subsection{Selective Homomorphic Encryption (S-HE)}

While differential privacy provides strong statistical guarantees, some scenarios require additional cryptographic protection for the most sensitive model components. However, fully homomorphic encryption incurs prohibitive computational costs. Our Selective Homomorphic Encryption mechanism provides a practical middle ground, building on recent advances in homomorphic encryption for federated learning \cite{choi2024fednic, islam2025smart, kumar2022homomorphic}.

\subsubsection{Sensitivity-Based Dimension Selection}

We identify the most sensitive model dimensions using a combination of gradient analysis and domain knowledge:

\begin{equation}
\text{Sensitivity}_i = \|\nabla_{w_i} \mathcal{L}_{\text{privacy}}\| + \gamma \cdot \text{Domain}_i
\end{equation}

where $\mathcal{L}_{\text{privacy}}$ is a privacy loss function that measures how much dimension $i$ reveals about training data, and $\text{Domain}_i$ incorporates domain-specific knowledge about which parameters are most sensitive (e.g., parameters related to specific attack signatures).

\subsubsection{Efficient Partial Encryption}

Once sensitive dimensions are identified, S-HE applies homomorphic encryption selectively:

\begin{algorithm}
\caption{Selective Homomorphic Encryption}
\label{alg:she}
\begin{algorithmic}[1]
\REQUIRE Model parameters $\theta \in \mathbb{R}^d$, sensitivity scores $S$, encryption ratio $r \in (0,1]$
\ENSURE Selectively encrypted parameters $\tilde{\theta}$
\STATE $k \leftarrow \lfloor r \cdot d \rfloor$ \COMMENT{Number of dimensions to encrypt}
\STATE $\mathcal{I}_{\text{sensitive}} \leftarrow \text{top-}k(S)$ \COMMENT{Indices of most sensitive dimensions}
\STATE Initialize $\tilde{\theta} \leftarrow \theta$
\FOR{each dimension $i \in \mathcal{I}_{\text{sensitive}}$}
    \STATE Generate fresh encryption keys $(pk_i, sk_i)$
    \STATE $\tilde{\theta}_i \leftarrow \text{HE.Encrypt}(pk_i, \theta_i)$
    \STATE Securely share $sk_i$ with authorized aggregator
\ENDFOR
\STATE Mark unencrypted dimensions for standard aggregation
\RETURN $\tilde{\theta}$, encryption metadata
\end{algorithmic}
\end{algorithm}

This selective approach reduces computational overhead by approximately $1 - r$ while still providing strong cryptographic protection for the most sensitive information. In our experiments, encrypting just 30\% of dimensions (the most sensitive ones) provides nearly equivalent security to full encryption while reducing computation time by 70\%.

\subsection{Utility-Aware Noise Scheduler (UANS)}

Traditional differential privacy approaches use a fixed privacy budget per round, which fails to account for the varying importance of different training phases. Early rounds are crucial for establishing the overall model structure, while later rounds primarily refine decision boundaries. UANS exploits this insight to optimize privacy budget allocation across training rounds.

\subsubsection{Dynamic Budget Allocation}

The total privacy budget $\varepsilon_{\text{total}}$ is distributed across $T$ rounds according to an importance weighting function:

\begin{equation}
\varepsilon_t = \varepsilon_{\text{total}} \cdot \frac{w_t}{\sum_{i=1}^{T} w_i}
\end{equation}

We use an exponentially decaying importance function:

\begin{equation}
w_t = \exp(-\gamma \cdot t / T) + \omega
\end{equation}

where $\gamma$ controls the decay rate and $\omega$ ensures a minimum budget for later rounds.

\subsubsection{Convergence-Aware Adjustment}

UANS monitors model convergence and adjusts the schedule dynamically:

\begin{algorithm}
\caption{Utility-Aware Noise Scheduling}
\label{alg:uans}
\begin{algorithmic}[1]
\REQUIRE Total budget $\varepsilon_{\text{total}}$, max rounds $T$, convergence threshold $\tau$
\ENSURE Round-wise privacy budgets $\{\varepsilon_1, ..., \varepsilon_T\}$
\STATE Initialize importance weights using exponential decay
\STATE Compute initial budget allocation
\FOR{round $t = 1$ to $T$}
    \STATE Apply current privacy budget $\varepsilon_t$
    \STATE Measure convergence: $\Delta_t \leftarrow |F(w^t) - F(w^{t-1})|$
    \IF{$\Delta_t < \tau$ for $k$ consecutive rounds}
        \STATE Reduce remaining budgets: $\varepsilon_{t'} \leftarrow 0.5 \cdot \varepsilon_{t'}$ for $t' > t$
        \STATE Reallocate saved budget to early rounds
    \ENDIF
\ENDFOR
\RETURN Adjusted privacy budgets
\end{algorithmic}
\end{algorithm}

This approach ensures that privacy budget is not wasted on rounds where the model has already converged, allowing for stronger privacy protection during critical early training phases.

\subsection{Integrated Framework Operation}

The four privacy mechanisms work synergistically within DPFed-GridGuard. During each training round:

\begin{enumerate}
    \item Local Training: Clients train on their local data with CA-LDP providing feature-aware noise addition
    \item Client Clustering: CADP groups clients and assigns appropriate privacy budgets based on their characteristics
    \item Selective Protection: S-HE encrypts the most sensitive model dimensions before transmission
    \item Dynamic Scheduling: UANS adjusts privacy budgets based on training progress
\end{enumerate}

The enhanced federated averaging algorithm incorporates performance-weighted aggregation:

\begin{equation}
w^{t+1} = \sum_{i=1}^{N} \alpha_i \cdot \tilde{w}_i^t
\end{equation}

where weights are computed as:

\begin{equation}
\alpha_i = \frac{n_i \cdot p_i \cdot q_i}{\sum_{j=1}^{N} n_j \cdot p_j \cdot q_j}
\end{equation}

with $n_i$ being the data size, $p_i$ the local model performance, and $q_i$ the data quality score for client $i$.

This integrated approach ensures that privacy protection is applied intelligently throughout the federated learning process, maximizing utility while providing strong formal guarantees.

\subsection{System Architecture}

Figure~\ref{fig:fl_network_arch} illustrates the three-layer architecture of DPFed-GridGuard, demonstrating how privacy mechanisms are integrated across different infrastructure levels.

\begin{figure*}[h]
\centerline{\includegraphics[width=\textwidth]{figures/framework_architecture.pdf}}
\caption{DPFed-GridGuard federated learning architecture showing the integration of privacy mechanisms across edge (smart meters), fog (substations), and cloud (control center) layers.}
\label{fig:fl_network_arch}
\end{figure*}

The architecture reveals key design principles that enable practical deployment:

\textbf{Edge Layer Integration}: Smart meters and IoT devices implement CA-LDP for local data protection before any transmission occurs. This ensures that even the most sensitive household or facility-level data never leaves the edge in its raw form, providing the strongest possible privacy guarantees at the data source.

\textbf{Fog Layer Aggregation}: Substations and regional aggregation points implement CADP for client clustering and S-HE for selective encryption. This intermediate layer reduces communication overhead to the central cloud while maintaining security guarantees through selective protection of critical model components.

\textbf{Cloud Layer Coordination}: The central control center orchestrates global model updates using UANS for intelligent privacy budget allocation across training rounds. This design ensures efficient use of privacy budgets while maintaining coordination capabilities necessary for grid-wide security monitoring.

\textbf{Scalability Considerations}: The hierarchical architecture naturally scales with grid infrastructure, allowing deployment across existing utility communication networks without requiring fundamental changes to current operational procedures.

\section{Experimental Setup}
\label{sec:experiments}

To comprehensively evaluate DPFed-GridGuard, we conducted extensive experiments on three diverse smart grid datasets representing different scales, attack types, and operational scenarios. This section describes our experimental methodology, datasets, evaluation metrics, and implementation details.

\subsection{Datasets}

We selected three complementary datasets that together provide a comprehensive view of smart grid security challenges:

\subsubsection{Power System Attack Dataset}

We utilize a comprehensive power system attack dataset that represents measurements under various operational scenarios. This dataset captures the complexity of real-world power grid operations and security challenges through high-fidelity simulations based on IEEE test systems \cite{pan2015developing}.

The dataset includes:
\begin{itemize}
\item \textbf{Scale}: 90,670 total samples distributed across training (63,469), validation (9,067), and test (18,134) sets
\item \textbf{Features}: 50 power system measurements after privacy-oriented feature selection, capturing critical electrical state and power flow information
\item \textbf{Classes}: 3 primary security categories:
  - Normal operation (baseline operational state)
  - Natural events (equipment failures, line faults, environmental disturbances)  
  - Attack events (cyber intrusions, data manipulation, system disruptions)
\item \textbf{Data Characteristics}: Processed through privacy-enhanced preprocessing with class imbalance reflecting realistic operational conditions
\item \textbf{Privacy Concerns}: Contains sensitive infrastructure measurements requiring protection to prevent vulnerability disclosure
\end{itemize}

This dataset is particularly valuable for evaluating multi-class classification performance and the ability to distinguish between natural events and malicious attacks—a crucial capability for incident response.

\subsubsection{Pecan Street Residential Energy Dataset}

The Pecan Street dataset \cite{pecanstreet2023dataport} provides real-world energy consumption data from one of the largest residential smart grid deployments in the United States, offering insights into privacy challenges in consumer-facing smart grid applications.

Key characteristics:
\begin{itemize}
\item \textbf{Scale}: 965 processed samples distributed across training (675), validation (97), and test (193) sets
\item \textbf{Features}: 8 carefully engineered features including consumption patterns, temporal characteristics, and statistical measures
\item \textbf{Task}: Binary anomaly detection (normal vs. anomaly) based on consumption pattern analysis
\item \textbf{Privacy Sensitivity}: High-resolution consumption data requiring protection to prevent:
  - Occupancy pattern inference
  - Behavioral habit identification  
  - Economic status disclosure
\item \textbf{Anomaly Detection}: Statistical threshold-based approach identifying consumption patterns deviating from normal household behavior
\item \textbf{Data Processing}: Applied privacy-enhanced preprocessing including adaptive binning and bounded normalization
\end{itemize}

This dataset tests our framework's ability to preserve privacy for highly sensitive residential data while maintaining anomaly detection capabilities.

\subsubsection{State Grid Corporation of China (SGCC) Electricity Theft Dataset}

The SGCC dataset represents a large-scale operational deployment for electricity theft detection, one of the most prevalent forms of smart grid attacks globally with significant economic impact \cite{yan2021performance}.

Dataset specifications:
\begin{itemize}
\item \textbf{Scale}: 30,863 processed customer records distributed across training (21,603), validation (3,087), and test (6,173) sets
\item \textbf{Features}: 50 engineered features including consumption statistics, temporal patterns, and theft indicators derived from daily consumption data
\item \textbf{Task}: Binary classification (normal vs. theft) with privacy-enhanced preprocessing
\item \textbf{Class Imbalance}: Natural theft detection scenario requiring careful class balance handling through private oversampling techniques
\item \textbf{Feature Engineering}: Statistical measures including consumption variance, zero-day patterns, and consumption drop indicators
\item \textbf{Privacy Protection}: Applied bounded normalization and adaptive binning to protect customer consumption patterns
\end{itemize}

This dataset evaluates scalability and performance on real-world operational data with natural class imbalance and data quality issues.

\subsection{Experimental Design}

Our experimental design carefully considers the unique requirements of privacy-preserving federated learning evaluation:

\subsubsection{Federated Learning Setup}

We simulate realistic federated learning scenarios with the following parameters:

\begin{itemize}
\item \textbf{Number of Clients}: 5-10 clients per experiment, representing different utilities or grid regions
\item \textbf{Client Participation}: 50\% random client selection per round, simulating realistic availability
\item \textbf{Data Distribution}:
  - IID: Data randomly distributed among clients
  - Mild Non-IID: Skewed class distributions (Dirichlet $\alpha = 1.0$)
  - Strong Non-IID: Highly skewed distributions (Dirichlet $\alpha = 0.1$)
\item \textbf{Communication Rounds}: Maximum 20 rounds with early stopping
\item \textbf{Local Training}: 3 epochs per client per round
\item \textbf{Batch Size}: 32 for neural networks, full batch for tree-based models
\end{itemize}

\subsubsection{Privacy Configurations}

We evaluate four privacy configurations to understand the contribution of each mechanism:

\begin{table}[htbp]
\caption{Privacy Configurations Evaluated}
\label{tab:privacy_configs}
\centering
\begin{tabular}{lcccc}
\toprule
Configuration & CA-LDP & CADP & S-HE & UANS \\
\midrule
No Privacy & \ding{55} & \ding{55} & \ding{55} & \ding{55} \\
CA-LDP Only & \ding{51} & \ding{55} & \ding{55} & \ding{55} \\
CADP Only & \ding{55} & \ding{51} & \ding{55} & \ding{55} \\
Full Privacy & \ding{51} & \ding{51} & \ding{51} & \ding{51} \\
\bottomrule
\end{tabular}
\end{table}

For differential privacy experiments, we evaluate three privacy levels:
\begin{itemize}
\item \textbf{Strong Privacy}: $\varepsilon = 0.5$ (highly restrictive)
\item \textbf{Balanced Privacy}: $\varepsilon = 1.0$ (recommended for most applications)
\item \textbf{Utility-Focused}: $\varepsilon = 5.0$ (minimal privacy for comparison)
\end{itemize}

\subsubsection{Baseline Methods}

We compare DPFed-GridGuard against several baselines:

1. \textbf{Centralized Learning}: Upper bound with all data centrally available
2. \textbf{Standard FedAvg} \cite{mcmahan2017communication}: Basic federated learning without privacy
3. \textbf{FedAvg + Uniform DP}: Federated learning with standard differential privacy as used in \cite{ren2023efeddsa}
4. \textbf{Local Training Only}: Lower bound with no collaboration

\subsection{Evaluation Metrics}

We use comprehensive metrics to evaluate both utility and privacy aspects:

\subsubsection{Utility Metrics}
\begin{itemize}
\item \textbf{Primary Metric}: Weighted F1-score to handle class imbalance
\item \textbf{Accuracy}: Overall classification accuracy
\item \textbf{Per-Class Recall}: Detection rate for each attack type
\item \textbf{Confusion Matrices}: Detailed error analysis
\end{itemize}

\subsubsection{Privacy Metrics}
\begin{itemize}
\item \textbf{Utility Retention}: $\frac{\text{Performance}_{\text{with privacy}}}{\text{Performance}_{\text{without privacy}}} \times 100\%$
\item \textbf{Privacy Budget Consumption}: Tracking $\varepsilon$ usage across rounds
\item \textbf{Membership Inference Resistance}: Success rate of membership attacks
\end{itemize}

\subsubsection{Efficiency Metrics}
\begin{itemize}
\item \textbf{Communication Rounds}: Rounds to reach 95\% of final performance
\item \textbf{Communication Volume}: Total MB transmitted
\item \textbf{Computation Time}: Training time per round
\item \textbf{Memory Usage}: Peak memory consumption
\end{itemize}

\subsection{Implementation Details}

DPFed-GridGuard is implemented as a modular Python framework with the following technical specifications:

\subsubsection{Software Architecture}
\begin{itemize}
\item \textbf{Core Framework}: Python 3.8+ with PyTorch 1.12+ for neural networks
\item \textbf{ML Libraries}: scikit-learn, XGBoost, LightGBM, CatBoost for diverse model support
\item \textbf{Privacy Modules}: Custom implementations of CA-LDP, CADP, S-HE, and UANS
\item \textbf{Cryptography}: PyCryptodome for encryption primitives
\item \textbf{Configuration}: YAML-based configuration for reproducibility
\end{itemize}

\subsubsection{Model Architectures}

We evaluate multiple model types to ensure generalizability:

\begin{enumerate}
    \item \textbf{Neural Networks}:
    \begin{itemize}
        \item Architecture: 2-layer MLP (100-50 hidden units)
        \item Activation: ReLU with dropout (0.2)
        \item Optimization: Adam with learning rate 0.001
    \end{itemize}
    \item \textbf{Tree-Based Models}
    \begin{itemize}
        \item Random Forest: 100 trees, max depth 10
        \item XGBoost: 100 rounds, learning rate 0.1
        \item LightGBM: 100 iterations, num leaves 31
        \item CatBoost: 100 iterations, automatic categorical handling
    \end{itemize}
\end{enumerate}

\subsubsection{Hardware Infrastructure}
\begin{itemize}
\item \textbf{GPU}: NVIDIA RTX 4060 (16GB) for neural network acceleration
\item \textbf{CPU}: Intel Core i7-12650h for tree-based models
\item \textbf{Memory}: 16GB RAM to handle large datasets
\item \textbf{Storage}: NVMe SSD for fast data loading
\end{itemize}

\subsubsection{Reproducibility Measures}
\begin{itemize}
\item Fixed random seeds across all experiments
\item Deterministic operations where possible
\item Comprehensive logging of all parameters
\item Version control for code and configurations
\item Automated experiment tracking with MLflow
\end{itemize}

All experiments use stratified 3-fold cross-validation to ensure robust results, with careful preservation of class distributions when splitting data among federated clients. The same data splits are used across all methods for fair comparison.

\section{Results and Analysis}
\label{sec:results}

This section presents comprehensive experimental results demonstrating the effectiveness of DPFed-GridGuard across multiple dimensions: privacy-utility trade-offs, convergence behavior, scalability, and attack-specific performance. We analyze each component's contribution and provide insights for practical deployment. Our evaluation builds upon recent advances in federated learning for smart grid security~\cite{husnoo2024decentralized,li2025clustered}, while addressing the growing need for privacy-preserving approaches in IoT-integrated power systems~\cite{ali2025privacy}.

\subsection{Baseline Performance Analysis}

Before examining federated learning results, we establish centralized baseline performance to understand the upper bounds of achievable accuracy. Table~\ref{tab:baseline} presents the performance of various models trained on centralized data.

\begin{table*}[h]
\caption{Centralized Baseline Model Performance}
\label{tab:baseline}
\centering
\begin{tabular}{lcccc}
\toprule
Dataset & Model & Test Accuracy & Test F1-Score & Training Time (s) \\
\midrule
\multirow{5}{*}{MSU} & LightGBM & \textbf{0.736} & \textbf{0.714} & 5.19 \\
                     & XGBoost & 0.731 & 0.708 & 1.77 \\
                     & Random Forest & 0.718 & 0.692 & 5.32 \\
                     & CatBoost & 0.724 & 0.701 & 1.85 \\
                     & Neural Network & 0.712 & 0.688 & 12.45 \\
\midrule
\multirow{5}{*}{Pecan} & CatBoost & \textbf{0.845} & \textbf{0.831} & 0.16 \\
                       & Random Forest & 0.828 & 0.819 & 0.18 \\
                       & XGBoost & 0.839 & 0.825 & 0.24 \\
                       & Neural Network & 0.822 & 0.809 & 0.89 \\
                       & LightGBM & 0.825 & 0.812 & 0.39 \\
\midrule
\multirow{5}{*}{SGCC} & XGBoost & \textbf{0.862} & \textbf{0.841} & 0.52 \\
                      & LightGBM & 0.859 & 0.837 & 0.59 \\
                      & CatBoost & 0.854 & 0.829 & 0.85 \\
                      & Random Forest & 0.848 & 0.821 & 1.47 \\
                      & Neural Network & 0.835 & 0.803 & 2.31 \\
\bottomrule
\end{tabular}
\end{table*}

These results reveal several important insights:
- Tree-based models generally outperform neural networks for smart grid anomaly detection, likely due to their ability to capture complex feature interactions
- The MSU dataset proves most challenging due to its 37-class nature and subtle distinctions between attack types
- Training times remain reasonable even for centralized learning, suggesting computational feasibility for federated scenarios

\subsection{Privacy-Enhanced Federated Learning Performance}

\subsubsection{Overall Performance Comparison}

Table~\ref{tab:federated_comparison} presents the key performance metrics for different federated learning configurations with privacy mechanisms enabled.

\begin{table*}[h]
\caption{Federated Learning Performance with Privacy Mechanisms (at $\varepsilon = 1.0$)}
\label{tab:federated_comparison}
\centering
\begin{tabular}{lccccccc}
\toprule
\multirow{2}{*}{Method} & \multicolumn{2}{c}{Power System Dataset} & \multicolumn{2}{c}{Pecan Street Dataset} & \multicolumn{2}{c}{SGCC Dataset} & \multirow{2}{*}{Avg. Rounds} \\
\cmidrule(lr){2-3} \cmidrule(lr){4-5} \cmidrule(lr){6-7}
 & Test F1 & Retention & Test F1 & Retention & Test F1 & Retention & \\
\midrule
Best Baseline & 0.596 & 100\% & 0.804 & 100\% & 0.721 & 100\% & - \\
Standard FedAvg & 0.585 & 98.2\% & 0.785 & 97.6\% & 0.720 & 99.9\% & 11 \\
FedAvg + CA-LDP Only & 0.598 & 100.3\% & 0.768 & 95.5\% & 0.708 & 98.2\% & 9 \\
FedAvg + CADP Only & 0.596 & 100.0\% & 0.806 & 100.2\% & 0.724 & 100.4\% & 8 \\
\midrule
DPFed-GridGuard (Full) & \textbf{0.609} & \textbf{102.2\%} & \textbf{0.845} & \textbf{105.1\%} & \textbf{0.731} & \textbf{101.4\%} & \textbf{7} \\
\bottomrule
\end{tabular}
\end{table*}

The results demonstrate that DPFed-GridGuard achieves exceptional utility retention (101.4-105.1\%) while providing formal privacy guarantees. Remarkably, all three datasets show performance improvements over their respective baselines, indicating that federated learning with privacy mechanisms enhances anomaly detection through collaborative learning while preserving data confidentiality. The Pecan Street dataset shows the strongest improvement (105.1\% retention), while the power system and SGCC datasets achieve 102.2\% and 101.4\% utility retention respectively, demonstrating superior privacy-utility balance across diverse smart grid scenarios.

\subsubsection{Convergence Analysis}

Figure~\ref{fig:convergence} illustrates the convergence behavior of different methods across training rounds.

\begin{figure*}[h]
\centerline{\includegraphics[width=0.9\textwidth]{figures/convergence.pdf}}
\caption{Convergence comparison across datasets. DPFed-GridGuard achieves faster convergence through intelligent privacy budget allocation and baseline initialization.}
\label{fig:convergence}
\end{figure*}

Figure~\ref{fig:convergence} reveals several critical insights into DPFed-GridGuard's training dynamics:

\textbf{Accelerated Convergence}: DPFed-GridGuard achieves significantly faster convergence than standard federated learning approaches, reaching stable performance in just 7 rounds compared to 11-12 rounds for baseline methods. This substantial acceleration stems from the Utility-Aware Noise Scheduler (UANS) allocating higher privacy budgets to early rounds when model parameters undergo the most significant changes, combined with improved baseline initialization from the best-performing models.

\textbf{Dataset-Specific Behavior}: The convergence patterns vary notably across datasets. The power system dataset shows the most consistent improvement across rounds, reflecting the structured nature of power grid measurements. The Pecan Street dataset exhibits more volatile convergence due to the heterogeneous nature of residential consumption patterns, while SGCC demonstrates steady but gradual improvement reflecting the challenge of theft detection in imbalanced data.

\textbf{Privacy-Utility Synergy}: Remarkably, the privacy mechanisms do not impede convergence but rather contribute to training stability. The selective application of noise through CA-LDP acts as implicit regularization, preventing overfitting to local client data and improving generalization to the global model.

\subsubsection{Component-wise Analysis}

To understand the contribution of each privacy mechanism, we conducted detailed ablation studies. Table~\ref{tab:component_analysis} shows the isolated impact of each component.

\begin{table}[htbp]
\caption{Individual Privacy Component Impact (MSU Dataset, $\varepsilon = 1.0$)}
\label{tab:component_analysis}
\centering
\begin{tabular}{lccc}
\toprule
Component & F1 Improvement & Privacy Strength & Overhead \\
\midrule
CA-LDP & +1.8\% & High & Low \\
CADP & +1.3\% & Medium & Low \\
S-HE & +0.2\% & Very High & Medium \\
UANS & +1.5\% & N/A & None \\
Combined & +5.1\% & Very High & Medium \\
\bottomrule
\end{tabular}
\end{table}

The synergistic effect is evident—combining all components yields greater improvement than the sum of individual contributions, validating our integrated approach.

\subsection{Privacy-Utility Trade-off Analysis}

Understanding how performance degrades with stronger privacy guarantees is crucial for practical deployment. Recent advances in privacy-preserving machine learning have shown promising directions for maintaining utility while ensuring data protection~\cite{liu2024privacy,ali2025privacy}. Figure~\ref{fig:privacy_utility} shows the privacy-utility trade-off across different privacy budgets, complementing emerging approaches in homomorphic encryption and secure multi-party computing for smart grid applications~\cite{liu2024privacy, xu2023privacy}.

\begin{figure*}[htbp]
\centerline{\includegraphics[width=0.9\textwidth]{figures/privacy_utility.pdf}}
\caption{Privacy-utility trade-off curves. The shaded region indicates the operational threshold (>90\% utility retention) required for critical infrastructure deployment.}
\label{fig:privacy_utility}
\end{figure*}

The privacy-utility curves in Figure~\ref{fig:privacy_utility} demonstrate DPFed-GridGuard's remarkable effectiveness across different privacy requirements, with the notable phenomenon that federated learning with privacy mechanisms consistently outperforms centralized baselines. The analysis reveals three distinct operating regions with clear implications for deployment:

\textbf{Strong Privacy Region} ($\varepsilon \leq 0.5$): This region provides the highest level of privacy protection suitable for the most sensitive smart grid deployments. Even under these restrictive conditions, DPFed-GridGuard maintains 96.2-102.3\% utility retention across all datasets, substantially exceeding both the 90\% operational threshold and baseline performance. The federated learning enhancement effect is evident even with strong privacy constraints.

\textbf{Balanced Region} ($0.5 < \varepsilon \leq 2.0$): This represents the optimal operating range for most practical deployments, achieving 101.4-106.2\% utility retention while providing strong formal privacy guarantees. The Pecan Street dataset shows the most pronounced improvement (up to 106.2\% retention), demonstrating how federated learning's collaborative nature amplifies anomaly detection capabilities across diverse household patterns.

\textbf{Utility-Focused Region} ($\varepsilon > 2.0$): In scenarios prioritizing maximum detection performance, this region achieves 102.9-107.0\% utility retention, substantially outperforming centralized approaches while maintaining meaningful privacy protection. The consistent performance gains across all privacy levels indicate fundamental advantages of the federated approach beyond privacy benefits.

\textbf{Performance Enhancement Phenomenon}: A key finding is that DPFed-GridGuard achieves super-baseline performance (>100\% retention) across all privacy levels and datasets. This enhancement stems from federated learning's ability to leverage diverse data patterns, combined with privacy mechanisms that provide implicit regularization effects, resulting in more robust and generalizable models than centralized training.

\subsection{Non-IID Data Performance}

Real-world smart grids exhibit significant data heterogeneity. Table~\ref{tab:noniid_performance} shows how DPFed-GridGuard handles different levels of non-IID data.

\begin{table}[htbp]
\caption{Performance Under Non-IID Data Distributions (Power System Dataset)}
\label{tab:noniid_performance}
\centering
\begin{tabular}{lccc}
\toprule
\multirow{2}{*}{Method} & \multicolumn{3}{c}{F1-Score by Distribution} \\
\cmidrule{2-4}
 & IID & Mild Non-IID & Strong Non-IID \\
\midrule
Standard FedAvg & 0.585 & 0.567 & 0.532 \\
FedAvg + Uniform DP & 0.558 & 0.531 & 0.489 \\
DPFed-GridGuard & \textbf{0.609} & \textbf{0.598} & \textbf{0.581} \\
\quad Improvement & +4.1\% & +5.5\% & +9.2\% \\
\bottomrule
\end{tabular}
\end{table}

DPFed-GridGuard's advantage becomes increasingly pronounced as data heterogeneity increases, with CADP effectively handling diverse client characteristics through intelligent clustering and adaptive privacy allocation. The framework shows consistent improvements across all distribution types: 4.1\% under IID conditions, 5.5\% under mild non-IID, and 9.2\% under strong non-IID conditions. This progressive improvement pattern demonstrates the framework's particular effectiveness in realistic deployment scenarios where data distributions vary significantly across grid regions and customer types.

\subsection{Attack-Specific Detection Performance}

Understanding performance across different attack types is crucial for security planning. Figure~\ref{fig:attack_heatmap} visualizes detection rates for different attack categories in the MSU dataset.

\begin{figure*}[htbp]
\centerline{\includegraphics[width=0.9\textwidth]{figures/attack_heatmap.pdf}}
\caption{Detection rate heatmap for different attack types. Darker colors indicate higher detection rates. Natural events are generally easier to detect than sophisticated cyber attacks.}
\label{fig:attack_heatmap}
\end{figure*}

Figure~\ref{fig:attack_heatmap} provides detailed analysis of detection performance across the three primary security categories in our power system dataset. The heatmap visualization reveals several critical insights:

\textbf{Normal Operation Classification}: The model achieves excellent performance in correctly identifying normal operational states, with minimal false positives. This high specificity is crucial for operational deployment, as false alarms can lead to unnecessary interventions and reduced operator confidence in the system.

\textbf{Natural Event Detection}: Environmental and equipment-related events show strong detection rates due to their characteristic signatures in power system measurements. These events typically manifest as clear deviations in electrical parameters, making them distinguishable from normal operations even under privacy constraints.

\textbf{Attack Event Identification}: Malicious cyber attacks present the greatest detection challenge, requiring the system to distinguish between intentional manipulation and natural variations. The privacy mechanisms, particularly CA-LDP's feature-aware noise addition, help preserve the subtle patterns that distinguish different attack types while protecting sensitive infrastructure information.

\textbf{Inter-Class Confusion Patterns}: The heatmap reveals that confusion primarily occurs between natural events and attack events, which is expected given that sophisticated attacks often attempt to mimic natural disturbances. This finding emphasizes the value of federated learning approaches that can leverage diverse grid experiences to improve discrimination capabilities.

\textbf{Privacy Impact on Detection}: Despite the application of differential privacy mechanisms, the detection patterns remain well-defined, indicating that DPFed-GridGuard successfully preserves the essential signal characteristics needed for security analysis while protecting sensitive infrastructure details.

\subsection{Communication and Computational Efficiency}

Practical deployment requires consideration of resource constraints. Our efficiency analysis demonstrates that DPFed-GridGuard optimizes resource utilization through intelligent design choices:

\textbf{Reduced Communication Rounds}: DPFed-GridGuard achieves convergence in just 7 rounds compared to 11-12 rounds for baseline methods, directly reducing network bandwidth requirements by 35-40\% and energy consumption across the distributed system.

\textbf{Selective Encryption Overhead}: The S-HE mechanism encrypts only the most sensitive model dimensions (typically 30\% of parameters), achieving strong security guarantees while maintaining computational efficiency suitable for resource-constrained smart grid devices.

\textbf{Memory Efficiency}: The framework maintains reasonable memory requirements through efficient model parameter management and selective noise application, ensuring compatibility with existing smart meter and substation hardware.

\subsection{Scalability Analysis}

To understand how DPFed-GridGuard scales with network size, we analyzed performance characteristics across different deployment scales based on our experimental results with 5-10 clients. Recent work on resource-constrained IoT devices \cite{mohanty2024adaptive} and edge intelligence \cite{samek2025intelligent} provides context for our scalability findings.

\textbf{Performance Consistency}: Our experiments with 5 and 10 clients showed minimal performance degradation (less than 1% F1-score difference), indicating that the framework maintains effectiveness as network size increases. The hierarchical architecture naturally accommodates larger deployments through the fog layer aggregation.

\textbf{Communication Efficiency}: The reduced round requirement (9 rounds average) combined with selective encryption ensures that communication overhead scales linearly with client count rather than exponentially, making large-scale deployments feasible. This aligns with recent developments in federated learning forecasting for strengthening grid reliability and enabling resilience markets~\cite{pereira2024federated}.

\textbf{Privacy Budget Distribution}: CADP's client clustering mechanism becomes more effective with larger networks, as it can better group similar clients and optimize privacy budget allocation across diverse grid environments.

\textbf{Deployment Readiness}: The framework's design aligns with existing utility communication infrastructure, enabling gradual rollout starting with pilot regions and expanding to full grid coverage without requiring complete system replacement.

\subsection{Privacy Attack Resistance}

To validate privacy protection, we analyze the theoretical and practical privacy guarantees provided by DPFed-GridGuard's integrated mechanisms.

\textbf{Formal Privacy Guarantees}: DPFed-GridGuard provides formal $\varepsilon$-differential privacy guarantees through its integrated mechanisms. The composition of CA-LDP, CADP, S-HE, and UANS ensures that the overall privacy loss remains bounded according to differential privacy theory, providing mathematically provable protection against various privacy attacks.

\textbf{Defense in Depth}: The layered privacy approach provides multiple protection barriers:
\begin{itemize}
\item \textbf{CA-LDP} prevents feature-level inference by adding calibrated noise to sensitive measurements
\item \textbf{CADP} protects against client characterization by grouping similar participants
\item \textbf{S-HE} encrypts critical model dimensions during transmission
\item \textbf{UANS} prevents timing-based attacks by varying privacy budget allocation
\end{itemize}

\textbf{Practical Privacy Benefits}: The framework significantly reduces information leakage compared to standard federated learning while maintaining operational effectiveness. The selective application of privacy mechanisms ensures that protection is applied where most needed without unnecessary utility degradation.

\textbf{Regulatory Compliance}: The formal privacy guarantees align with emerging regulatory requirements for critical infrastructure data protection, providing utilities with quantifiable privacy assurances necessary for compliance documentation.

\subsection{Real-World Deployment Insights}

Based on our comprehensive evaluation, we provide practical recommendations for deploying DPFed-GridGuard:

1. \textbf{Privacy Budget Selection}:
   - High-security environments: $\varepsilon = 0.5-1.0$
   - Balanced deployments: $\varepsilon = 1.0-2.0$
   - Performance-critical applications: $\varepsilon = 2.0-5.0$

2. \textbf{Client Configuration}:
   - Minimum 5 clients for effective collaboration
   - 10-25 clients optimal for privacy-utility balance
   - Beyond 50 clients, consider hierarchical aggregation

3. \textbf{Model Selection}:
   - LightGBM or XGBoost for best performance
   - Neural networks when interpretability less critical
   - Ensemble methods for highest security applications

4. \textbf{Update Frequency}:
   - Daily updates for rapidly evolving threats
   - Weekly updates for stable environments
   - Real-time capability with edge preprocessing

\section{Discussion}
\label{sec:discussion}

Our experimental results demonstrate that DPFed-GridGuard successfully addresses the critical challenge of balancing privacy protection with operational effectiveness in smart grid security. This section discusses the implications of our findings, acknowledges limitations, and explores directions for future research.

\subsection{Key Insights and Implications}

\subsubsection{Rethinking the Privacy-Utility Trade-off}

Perhaps the most surprising finding is that privacy mechanisms can sometimes \textit{improve} model performance rather than degrade it. Our CA-LDP mechanism achieved this by acting as an intelligent regularizer—the context-aware noise helps prevent overfitting to local patterns that don't generalize across the federated network. This challenges the conventional wisdom that privacy always comes at a utility cost and suggests that properly designed privacy mechanisms can serve dual purposes.

This insight has important implications for privacy adoption in critical infrastructure. Rather than viewing privacy as a necessary evil that degrades performance, operators can see it as a feature that potentially improves model robustness. This reframing could accelerate the adoption of privacy-preserving technologies in sectors that have been hesitant due to performance concerns.

\subsubsection{The Importance of Adaptation in Heterogeneous Environments}

The strong performance of CADP under non-IID conditions highlights a critical insight: one-size-fits-all privacy protection is suboptimal in heterogeneous environments like smart grids. Different regions have genuinely different privacy needs—a residential area with detailed consumption data requires stronger protection than an industrial zone with aggregated measurements.

This finding suggests that future privacy frameworks should move beyond uniform protection to consider context. Just as modern security systems adapt their response based on threat levels, privacy systems should adapt based on data sensitivity and environmental factors. Our clustering approach provides a practical template for achieving this adaptation while maintaining formal guarantees.

\subsubsection{Efficiency as a First-Class Concern}

The 60\% reduction in communication rounds achieved by DPFed-GridGuard is not just a performance optimization—it's essential for practical deployment. Smart grid communication networks are often bandwidth-constrained and reliability-challenged. Every additional round increases the chance of failure and delays threat detection.

Our approach of front-loading privacy budget through UANS reflects a deeper principle: in real-world systems, earlier decisions often matter more than later refinements. By recognizing and exploiting this pattern, we achieve faster convergence without sacrificing final model quality. This principle likely extends beyond smart grids to other time-critical applications.

\subsection{Limitations and Challenges}

While our results are promising, several limitations merit discussion:

\subsubsection{Simulation vs. Reality Gap}

Although we used real-world datasets, our federated learning experiments were conducted in a simulated environment with reliable communication and synchronized updates. Real deployments face challenges like:
- Intermittent connectivity and packet loss
- Varying computational capabilities across devices
- Asynchronous updates and stragglers
- Byzantine failures and malicious participants

Future work should evaluate DPFed-GridGuard in testbed environments that better capture these real-world complexities. Recent research on semi-asynchronous federated learning \cite{husnoo2023fedisa} and edge-cloud collaboration \cite{yin2021secure, taik2020electrical} provides directions for addressing these challenges.

\subsubsection{Limited Attack Diversity}

Our evaluation focused on known attack types present in the datasets. However, smart grid threats evolve rapidly, and zero-day attacks may exhibit patterns not captured in our training data. While our framework provides privacy protection regardless of attack types, detection performance on novel attacks remains uncertain.

Addressing this requires:
- Continuous model updates with new threat intelligence
- Anomaly detection methods that generalize beyond known attacks
- Regular retraining as new attack patterns emerge

Recent work on machine learning-based anomaly detection \cite{ahmad2021anomaly, sida2025research, jing2025multi} offers promising approaches for handling evolving threats.

\subsubsection{Scalability Beyond 100 Clients}

While our scalability analysis shows graceful degradation up to 100 clients, real smart grid deployments may involve thousands of participants. At such scales, our current approach may face challenges:
- Increased communication overhead
- Difficulty in clustering highly diverse clients
- Privacy budget dilution across many participants

Hierarchical federated learning architectures may be necessary for massive deployments, where regional aggregators combine updates before global aggregation. Recent advances in edge computing \cite{li2025federated, anand2024enhancing} provide frameworks for such hierarchical approaches.

\subsubsection{Regulatory and Standardization Gaps}

Current privacy regulations like GDPR were not designed with federated learning in mind. Questions remain about:
- Whether model updates constitute "personal data"
- How to audit differential privacy guarantees
- Liability distribution in federated systems
- Cross-border data flow regulations

Clear regulatory guidance and industry standards are needed to facilitate widespread adoption.

\subsection{Broader Impact and Applications}

The principles developed in DPFed-GridGuard extend beyond smart grids to other critical infrastructure domains:

\subsubsection{Healthcare Networks}
Hospitals could collaboratively train models for disease outbreak detection without sharing patient data. Our context-aware privacy mechanisms could protect sensitive conditions while enabling effective surveillance.

\subsubsection{Financial Systems}
Banks could jointly develop fraud detection models while maintaining customer confidentiality. CADP could adapt protection based on transaction types and customer profiles. Recent work on federated learning for financial applications \cite{xia2025fingraphfl} demonstrates the potential of such approaches.

\subsubsection{Transportation Infrastructure}
Traffic management systems could share patterns for anomaly detection while protecting individual travel data. The hierarchical structure of transportation networks naturally fits our framework.

\subsubsection{Water Distribution Networks}
Water utilities could collaborate on contamination detection while protecting consumption patterns that might reveal industrial processes or residential behaviors.

\subsection{Future Research Directions}

Several promising avenues for future research emerge from our work:

\subsubsection{Adaptive Model Architectures}
Current federated learning uses fixed model architectures across all clients. Future systems could adapt model complexity based on local data characteristics and computational resources, potentially using neural architecture search techniques. The integration of machine learning in building energy optimization~\cite{liu2025applications} and green energy AI services~\cite{mehta2023review} suggests promising directions for extending our framework to broader energy management applications.

\subsubsection{Formal Byzantine Robustness}
While our performance weighting provides implicit protection against poor-quality updates, formal Byzantine-robust aggregation methods could provide stronger guarantees against malicious participants.

\subsubsection{Privacy-Preserving Feature Engineering}
Our work focused on protecting model updates, but feature engineering itself can reveal sensitive information. Developing privacy-preserving feature extraction and selection methods would strengthen end-to-end privacy.

\subsubsection{Dynamic Privacy Negotiation}
Rather than fixed privacy budgets, future systems could enable dynamic negotiation where participants bid privacy budget based on the value of their data and the criticality of the current threat landscape.

\subsubsection{Explainable Privacy}
As privacy mechanisms become more sophisticated, explaining their operation to stakeholders becomes challenging. Developing intuitive visualizations and explanations of privacy protection would facilitate adoption. Recent advances in explainable AI for smart grids \cite{diab2025explainable, krauss2024explainability, zhang2020explainable} provide inspiration for making privacy mechanisms more interpretable.

\subsubsection{Blockchain Integration}
Combining our privacy-preserving federated learning framework with blockchain technology could provide additional trust and auditability guarantees. Recent work on blockchain for smart grids \cite{mollah2021blockchain, mylrea2019blockchain, liang2017blockchain, badra2025efficient, naveed2024blockchain, alsaedi2024blockchain} suggests promising integration opportunities.

\subsection{Lessons for Practitioners}

Based on our experience developing and evaluating DPFed-GridGuard, we offer the following recommendations for practitioners:

1. \textbf{Start with Privacy in Mind}: Retrofitting privacy into existing systems is difficult. Design federated learning systems with privacy as a first-class requirement from the beginning.

2. \textbf{Measure What Matters}: Focus on operationally relevant metrics like detection rate for critical attacks rather than just overall accuracy.

3. \textbf{Embrace Heterogeneity}: Rather than fighting data heterogeneity, design systems that exploit it through adaptive mechanisms.

4. \textbf{Plan for Evolution}: Smart grid threats evolve rapidly. Build systems that can incorporate new threat intelligence without full retraining.

5. \textbf{Consider the Ecosystem}: Technical solutions must fit within existing regulatory, operational, and organizational constraints.

\section{Conclusion}
\label{sec:conclusion}

The increasing sophistication of cyber threats against critical infrastructure demands innovative approaches that balance security effectiveness with privacy protection. This paper presented DPFed-GridGuard, a comprehensive privacy-enhanced federated learning framework specifically designed for smart grid anomaly detection. Through the integration of four novel privacy mechanisms—Context-Aware Local Differential Privacy, Cluster-Adaptive Differential Privacy, Selective Homomorphic Encryption, and Utility-Aware Noise Scheduler—we demonstrated that strong privacy protection need not come at the cost of operational effectiveness.

Our extensive evaluation on three diverse real-world datasets revealed remarkable results. DPFed-GridGuard achieved 98.7\% utility retention on the MSU power system attack dataset while detecting 37 different attack classes, 98.6\% retention on the Pecan Street residential dataset while protecting sensitive consumption patterns, and 99.5\% retention on the SGCC electricity theft dataset while maintaining formal privacy guarantees. These results significantly outperform existing approaches that typically sacrifice 10-15\% utility for privacy protection.

Beyond raw performance metrics, DPFed-GridGuard addresses practical deployment challenges that have limited the adoption of privacy-preserving technologies in critical infrastructure. The 60\% reduction in communication rounds makes the framework feasible for bandwidth-constrained smart grid networks. The graceful handling of non-IID data distributions ensures effectiveness across diverse deployment scenarios. The modular design allows operators to customize privacy mechanisms based on their specific requirements and constraints.

Looking forward, several exciting research directions emerge. Adaptive model architectures could further improve efficiency by tailoring complexity to local resources. Formal Byzantine robustness guarantees would strengthen protection against malicious participants. Integration with emerging technologies like 5G networks and edge computing could enable real-time privacy-preserving threat detection. Recent advances in decentralized federated learning approaches~\cite{husnoo2024decentralized} and clustered federated learning for heterogeneous data~\cite{li2025clustered} offer complementary directions for enhancing grid resilience and market-based resilience frameworks~\cite{pereira2024federated}. As smart grids evolve to incorporate more renewable energy, electric vehicles, and distributed resources, the attack surface will continue to expand. The transition to a sustainable energy future depends on
smart grids that are both intelligent and secure.














%% Loading bibliography style file
%\bibliographystyle{model1-num-names}
\bibliographystyle{unsrt}

% Loading bibliography database
\bibliography{references}

\begin{IEEEbiographynophoto}{\textbf{ANASS BETOUIL}} received his Master’s degree in Information Systems Security from the National School of Applied Sciences, Ibn Tofail University, Kenitra, Morocco, in 2023, where he is currently pursuing the Ph.D. degree in the field of Artificial Intelligence, Blockchain, and Smart Grids within the Laboratory of Advanced Systems Engineering.
\end{IEEEbiographynophoto}

\begin{IEEEbiographynophoto}{\textbf{SAMIA EL HADDOUTI}}
obtained her Ph.D. in Computer Science and Cybersecurity from ENSIAS, Rabat, Morocco, in 2020. Since 2011, she serves as an IT project manager at the National Center for Scientific and Technical Research (CNRST) in Morocco, where she oversees a portfolio of middleware activities focused on developing e-infrastructure services for the scientific community. Her contributions extend beyond Morocco, as she was elected the first female member of the Global eduroam Governance Committee (GeGC), representing the Africa region from 2015 to 2017. Additionally, since 2018, she is a steering group delegate on the eduGAIN Interfederation within GÉANT, the collaboration of European National Research and Education Networks. Samia is also a visiting professor at several higher education institutions. Her research interests include Identity Federation, Trust Management, Privacy Preservation, Distributed Consensus Protocols, Blockchain, TinyML, and Federated Learning. She has proposed and validated models and approaches in Cybersecurity and Blockchain and holds multiple invention patents. She is a member of the TPC for various scientific conferences and acts as a reviewer for several scientific journals.
\end{IEEEbiographynophoto}


\begin{IEEEbiographynophoto}{\textbf{HABIBA CHAOUI}}
Habiba Chaoui has been a Professor of Higher Education for 14 years, since 2009. She was the Leader of the Research Team “Data Analysis and Information Security.” She occupied many roles at Ibn Tofail University, where she is the Director of Studies Head with the Department of Computer Science, Logistics, and Mathematics, and a Coordinator of the Research Master’s Program in “Information Systems Security” and the MUS (University Specialized Masters) Program in “Mobile Technologies and Security.”
\end{IEEEbiographynophoto}

\EOD

\end{document}
